\Hefteintrag{1.3}{Was ist die DSGVO?}{

Die \emphColA{europäische Datenschutzgrundverordnung (DSGVO)} ist eine Verordnung der Europäischen Union, die die Regeln zur Verarbeitung personenbezogener Daten vereinheitlicht. Seit Mai 2018 ist sie in der EU sowie Island, Liechtenstein und Norwegen geltendes Recht. 

Die DSGVO besteht aus 99 Artikeln in 11 Kapiteln. Eine gravierende Neuregelung der DSGVO (neben einigen weiteren) war, dass \emphColB{Bußgelder} für Verstöße \emphColB{drastisch erhöht} wurde und die \emphColC{Verarbeitung personenbezogener Daten nur} aufgrund eines \emphColC{Erlaubnistatbestands} (Artikel 6) \emphColC{zulässig} ist:
\begin{itemize}
    \item Die betroffene Person hat ihre Einwilligung gegeben \emph{(das ist häufig in AGBs o.ä. versteckt)} 
    \item die Verarbeitung ist für die Erfüllung eines Vertrags oder zur Durchführung vorvertraglicher Maßnahmen erforderlich (z.B. \LoesungLuecke{Adresse für Versand}{6cm}).
    \item die Verarbeitung ist zur Erfüllung einer rechtlichen Verpflichtung erforderlich (z.B. \LoesungLuecke{Handel mit waffenfähigen Chemikalien}{15cm}).
    \item die Verarbeitung ist erforderlich, um lebenswichtige Interessen zu schützen (z.B. \LoesungLuecke{bei Entführungen}{15cm}).
    \item die Verarbeitung ist für die Wahrnehmung einer Aufgabe erforderlich, die im öffentlichen Interesse liegt (z.B. \LoesungLuecke{für Ermittlungen zur Gefahrenabwehr}{15cm}).
    \item die Verarbeitung ist zur Wahrung der berechtigten Interessen des Verantwortlichen oder eines Dritten erforderlich (z.B. \LoesungLuecke{Schufa-Auskunft bei Vermietung}{15cm}).
\end{itemize}

}